\subsection{Sammensatte funktioner}
Vi ved, at vi kan indsætte tal i stedet for $x$ i en funktion, f.eks.:
\begin{align*}
    f(x) &= 3x + 1 \\
    f(0) &= 3 \cdot 0 + 1 = 1 \\
    f(1) &= 3 \cdot 1 + 1 = 4 \\
    \vdots
\end{align*}

Vi kan også indsætte en funktion i en anden funktion. Lad os antage, at vi har to funktioner:
\begin{equation*}
    f(x) = 3x + 1 \qquad g(x) = 0.5x
\end{equation*}

Den sammensatte funktion $f(g(x))$ betyder, at vi indsætter udtrykket for $g(x)$ i stedet for $x$ i $f(x)$:
\begin{equation*}
    f(g(x)) = 3g(x) + 1
\end{equation*}

Indsætter vi $g(x) = 0.5x$, får vi:
\begin{align*}
    f(g(x)) &= 3 \cdot 0.5x + 1 \\
            &= 1.5x + 1
\end{align*}

\subsection{Inverse funktioner}
En invers funktion er en funktion, som ophæver en anden funktion. Når en funktion sammensættes med sin inverse, fås identitetsfunktionen:
\begin{equation*}
    f(f^{-1}(x)) = x \qquad \text{og} \qquad f^{-1}(f(x)) = x
\end{equation*}

Eksempel:
\begin{equation*}
    f(x) = 2x \qquad g(x) = 0.5x
\end{equation*}

Her gælder:
\begin{equation*}
    f(g(x)) = 2 \cdot 0.5x = x
\end{equation*}

Dermed er $g(x)$ den inverse funktion til $f(x)$, og vi kan skrive:
\begin{equation*}
    f^{-1}(x) = 0.5x
\end{equation*}

Man kender også inverse funktioner fra trigonometriske funktioner, f.eks.:
\begin{equation*}
    \cos^{-1}(\cos(x)) = x
\end{equation*}
(bemærk dog, at dette kun gælder for bestemte intervaller).

\subsubsection*{Bestemmelse af en invers funktion}
Vi kan finde den inverse funktion ved at isolere $x$.

Eksempel:
\begin{equation*}
    h(x) = 1.5x + 1
\end{equation*}

Sæt $y = h(x)$ og isolér $x$:
\begin{align*}
    y &= 1.5x + 1 \\
    y - 1 &= 1.5x \\
    x &= \frac{y - 1}{1.5}
\end{align*}

Den inverse funktion fås ved at bytte $y$ ud med $x$:
\begin{equation*}
    h^{-1}(x) = \frac{x - 1}{1.5}
\end{equation*}

(Man kan også skrive dette som $h^{-1}(x) = \frac{2}{3}x - \frac{2}{3}$).

\subsubsection*{Kontrol}
Vi kan kontrollere resultatet:
\begin{align*}
    h^{-1}(h(x)) &= \frac{1.5x + 1 - 1}{1.5} \\
                 &= \frac{1.5x}{1.5} \\
                 &= x
\end{align*}

Dermed er den inverse funktion korrekt.